\section{Introduction}\label{sec:introduction}
Reinforcement learning (RL) lets a robot learn navigation by trial and error, but training on a
real robot is slow, costly and unsafe \cite{sutton2018rl, zhao2020sim}. Policies are therefore
trained in simulation, where they often fail once moved to the real robot because the simulator
differs from reality. This is known as the sim-to-real gap \cite{zhao2020sim,
peng2018dynamics}. It can be narrowed in two ways: by making the simulator a closer copy of
reality, or by making the policy robust to the differences that remain, for example with domain
randomisation \cite{tobin2017domain, peng2018dynamics}.

This report tests the first approach on the Pepper robot, while applying the second equally in both
simulators, whose hardware is too fragile for training on the robot itself and for which no open
simulator yet supports RL navigation \cite{sekkat2024pepper}. A real room was scanned with a tablet
and turned into a photogrammetric simulator. A second, standard simulator of the same room was built
by hand from tape measurements. Navigation policies were trained with Soft Actor-Critic (SAC) under
the same domain randomisation in each simulator and then deployed, without further training, on the
physical Pepper over the same 20 routes. Because the source of the room's geometry is the only
difference between the two, the comparison isolates the effect of the reconstruction, including its
scanning artefacts.

Using odometry-estimated success on the real robot, the photogrammetry policy showed a smaller
performance drop from simulation to reality than the standard policy, and it also reached the goal
more often, although that difference was not statistically significant. Real-robot failures for both
policies happened at the same two obstacles, where the robot's built-in collision protection stopped
it.
